% !TeX root = ../drafts/03-proving-primitives.tex
\section{Proving primitives}\label{sec:prelim}

\begin{definition}[Iverson bracket]\label{def:iverson}
For a proposition $P$, \ $[P]=1$ if $P$ holds and $[P]=0$ otherwise.
\end{definition}

\begin{definition}[Multiset]\label{def:multiset}
A \emph{multiset} $\mset{\cdots}$ over a set $S$ is a multiplicity function $S\to\mathbb{N}$; $\uplus$ adds multiplicities.
\end{definition}

\begin{definition}[Multilinear extension]\label{def:mle}
A \emph{multilinear} is given by its $2^{\kappa}$ values on the boolean hypercube, a function $f:\cube{\kappa}\to\E$. Its \emph{multilinear extension} $\mle{f}:\E^{\kappa}\to\E$ is the unique polynomial of degree at most one in each variable agreeing with $f$ on $\cube{\kappa}$; we identify the two freely.
\end{definition}


\begin{definition}[Equality indicator]\label{def:eq}
For $X,Y\in\E^{\kappa}$,
\[
  \eq(X,Y)= \prod_i\bigl(X_iY_i+(1+X_i)(1+Y_i)\bigr) = \prod_{i=1}^{\kappa}\bigl(1+X_i+Y_i\bigr).
\]
\end{definition}

$\eq$ is multilinear, and $\eq(X,Y)=[X=Y]$ on the boolean hypercube $\cube{2\kappa}$.

\begin{definition}[Inner product over the boolean hypercube]\label{def:inner}
For $a,b$ on $\cube{\kappa}$, \ $\inner{a}{b}=\sum_{x\in\cube{\kappa}}a(x)\,b(x)$.
\end{definition}

\begin{fact}[$\eq$ carries the interpolation weights]\label{fact:eq}
For an arbitrary point $r\in\E^{\kappa}$,
\[
  \mle{f}(r)=\sum_{x\in\cube{\kappa}}f(x)\,\eq(r,x)=\inner{f}{\eq(r,\cdot)}.
\]
Binding only the first $\ell\le\kappa$ variables gives the same identity with the rest left free: for $r\in\E^{\ell}$ and $x\in\cube{\kappa-\ell}$,
\[
  \mle{f}(r,x)=\sum_{u\in\cube{\ell}}\eq(r,u)\,f(u,x),
\]
so a partial evaluation is again a multilinear, its values the $\eq(r,\cdot)$-weighted combination of the rows $f(u,\cdot)$.
\end{fact}


\begin{definition}[Table, column]\label{def:column}
A \emph{column} of height $2^{\kappa}$ is a function $\cube{\kappa}\to\K$ committed at the beginning of the proof (via \S\ref{sec:stacking}), and a \emph{table} is a group of columns, of equal height, connected by polynomial constraints. A $192$-bit \emph{value} (such as memory words) is indirectly committed by three columns, read as $V=V_0+V_1y+V_2y^{2}$.
\end{definition}

\begin{lemma}[Every element outside $\K$ is a cubic generator]\label{lem:kbasis}
For $\alpha\in\E\setminus\K$, the set $\{1,\alpha,\alpha^{2}\}$ is a $\K$-basis of $\E$.
\end{lemma}

\begin{proof}
Degrees multiply in the tower $\K\subseteq\K(\alpha)\subseteq\E$, so $[\E:\K(\alpha)]\,[\K(\alpha):\K]=[\E:\K]=3$. That degree is prime and $[\K(\alpha):\K]>1$ because $\alpha\notin\K$, hence $[\K(\alpha):\K]=3$ and $\K(\alpha)=\E$. The minimal polynomial of $\alpha$ over $\K$ then has degree $3$, so $1,\alpha,\alpha^{2}$ are $\K$-independent, and three independent elements of a space of $\K$-dimension $3$ are a basis.
\end{proof}

\begin{lemma}[Schwartz--Zippel]\label{lem:sz}
A nonzero $P\in\E[X_1,\dots,X_{\kappa}]$ of total degree at most $d$ satisfies
\[
  \Pr_{r\leftarrow\E^{\kappa}}\bigl[P(r)=0\bigr]\;\le\;\frac{d}{|\E|}.
\]
A nonzero multilinear has total degree at most $\kappa$, so it vanishes at a random point with probability at most $\kappa/|\E|$.
\end{lemma}

\begin{proof}
Induction on $\kappa$. For $\kappa=1$ a nonzero univariate of degree at most $d$ has at most $d$ roots, so a uniform $r_1$ is one of them with probability at most $d/|\E|$. For $\kappa>1$ write $P$ in its last variable,
\[
  P(X_1,\dots,X_{\kappa})=\sum_{k=0}^{d}X_{\kappa}^{\,k}\,Q_k(X_1,\dots,X_{\kappa-1}),
\]
and let $k^{\ast}$ be the largest $k$ with $Q_k\neq0$. Then $Q_{k^{\ast}}$ is nonzero in $\kappa-1$ variables of total degree at most $d-k^{\ast}$, so by induction $Q_{k^{\ast}}(r_{<\kappa})=0$ with probability at most $(d-k^{\ast})/|\E|$. Fix $r_{<\kappa}$ with $Q_{k^{\ast}}(r_{<\kappa})\neq0$: then $X_{\kappa}\mapsto P(r_{<\kappa},X_{\kappa})$ has degree exactly $k^{\ast}$, hence at most $k^{\ast}$ roots, and a uniform $r_{\kappa}$ is one with probability at most $k^{\ast}/|\E|$. Adding the two bounds gives $d/|\E|$.
\end{proof}

\begin{corollary}[Testing an identity at one random point]\label{cor:idtest}
Let $A,B\in\E[X_1,\dots,X_{\kappa}]$ have total degree at most $d$, and let $r\leftarrow\E^{\kappa}$ be drawn after both are fixed. If $A\neq B$ then
\[
  \Pr_{r}\bigl[A(r)=B(r)\bigr]\;\le\;\frac{d}{|\E|}.
\]
\end{corollary}

\begin{proof}
Apply Lemma~\ref{lem:sz} to $A-B$, nonzero of total degree at most $d$.
\end{proof}


\begin{fact}[Sumcheck~\cite{LFKN92}]\label{fact:sumcheck}
Over $\kappa$ rounds, sumcheck reduces a claim $\sum_{x\in\cube{\kappa}}P(x)=s$, on a $P$ of degree at most $d$ in each variable, to one evaluation $P(\fc)$ at a random $\fc\in\E^{\kappa}$. Round $j$ carries a claim $s_{j-1}$, with $s_0=s$: the prover sends the univariate
\[
  h_j(Y)=\sum_{x\in\cube{\kappa-j}}P(\fc_1,\dots,\fc_{j-1},Y,x)
\]
of degree at most $d$, the verifier checks $h_j(0)+h_j(1)\qeq s_{j-1}$, samples $\fc_j\leftarrow\E$ and takes $s_j=h_j(\fc_j)$. It ends holding $s_{\kappa}=P(\fc)$, for the caller to settle. A false $s$ survives with probability at most $d\kappa/|\E|$.
\end{fact}

\begin{definition}[Zerocheck]\label{def:zerocheck}
To prove that $P$ vanishes on $\cube{\kappa}$, the verifier samples $r\in\E^{\kappa}$ and the two parties run sumcheck on
\[
  \sum_{x\in\cube{\kappa}}P(x)\,\eq(r,x)\qeq 0 .
\]
By Fact~\ref{fact:eq} that sum is $\mle{g}(r)$, for $g$ the restriction of $P$ to the boolean hypercube: the zero polynomial if $P$ vanishes there, and otherwise a nonzero multilinear, which a random $r$ sends to $0$ with probability at most $\kappa/|\E|$ (Lemma~\ref{lem:sz}).
\end{definition}

\begin{lemma}[Partially fixed zerocheck~\cite{DT24}]\label{lem:fixed-zerocheck}
Let $g:\cube{\ell}\times\cube{m}\to\Ftwo$ be a nonzero Boolean multilinear. Fix $a\in\E^{\ell}$ such that the values $\{\eq(a,b):b\in\cube{\ell}\}$ are linearly independent over $\Ftwo$. If $s\leftarrow\E^m$ is sampled uniformly after $g$ is fixed, then
\[
  \Pr_s\bigl[\mle{g}(a,s)=0\bigr]\leq\frac{m}{|\E|}.
\]
\end{lemma}

\begin{proof}
Partially evaluate the first $\ell$ variables at $a$. By Fact~\ref{fact:eq}, the resulting multilinear $h:\cube{m}\to\E$ has values
\[
  h(y)=\mle{g}(a,y)=\sum_{b\in\cube{\ell}}g(b,y)\eq(a,b).
\]
Because $g$ is nonzero, there is some $y\in\cube{m}$ for which the coefficients $\{g(b,y):b\in\cube{\ell}\}$ are not all zero. These coefficients lie in $\Ftwo$, so the assumed $\Ftwo$-linear independence of the values $\eq(a,b)$ gives $h(y)\neq0$. Thus $h$ is a nonzero multilinear. Its extension satisfies $\mle{h}(s)=\mle{g}(a,s)$, and Lemma~\ref{lem:sz} gives the claimed probability bound.
\end{proof}

\paragraph{Two savings on a round message~\cite{Gru24}.} A round polynomial of degree $D$ travels as its coefficients, and would take $D+1$ of them, but the round check $h_j(0)+h_j(1)\qeq s_{j-1}$ already pins one to the incoming claim, so $D$ are enough. In coefficients that check is an addition, $\sum_{i\ge1}c_i\qeq s_{j-1}$, pinning $c_1$; the verifier then reads $h_j$ at its challenge by Horner, with no interpolation. The same idea applies to the known factor: a zerocheck's summand carries $\eq(r,\cdot)$, which the verifier can form itself, so only the cofactor travels and it is one degree lower. A zerocheck round on a degree-$d$ cofactor therefore costs $d$ field elements.

\begin{definition}[Inner-product commitment scheme]\label{def:ipcs}
A \emph{weighted claim} about a multilinear $g:\cube{M}\to\K$ is a pair $(W,c)$ asserting $\inner{W}{g}=c$, with $c\in\E$ and the weight $W:\cube{M}\to\E$ \emph{MLE-friendly}: the verifier can evaluate $\mle{W}$ anywhere in $\mathrm{poly}(M)$ operations. An \emph{inner-product commitment scheme} commits to $g$ and later proves any such claim,
\[
  \inner{W}{g}=\sum_{w\in\cube{M}}W(w)\,g(w)=c .
\]
The commitment is over $\K$, the opening over $\E$. An evaluation claim $\mle{g}(r)=c$ is the case $W=\eq(r,\cdot)$, by Fact~\ref{fact:eq}.
\end{definition}

We use WHIR~\cite{ACFY25} / Ligerito~\cite{NA25} as inner-product commitment scheme (see annex \ref{annex:pcs}). The two are essentially the same scheme when WHIR is instantiated with an interleaved code and Ligerito with a Reed--Solomon code.
